%%%%%%%%%%%%%%%%%%%%%%%%%%%%%%%%%%%%%%%%%%%%%%%%%%%%%%%%%%%%%%%%%%%%%%%%%%%%%%%%
\section{Etapa de calibración}
\label{sec:exp_calibracion}
%%%%%%%%%%%%%%%%%%%%%%%%%%%%%%%%%%%%%%%%%%%%%%%%%%%%%%%%%%%%%%%%%%%%%%%%%%%%%%%%

\textcolor{red}{%
[SCAFFOLDING --- Esta sección describe la etapa de calibración, cuyo objetivo fue determinar los valores óptimos de los hiperparámetros $\lambda$ (peso del costo dinámico) y $\mu$ (peso del costo de trayectoria) de la función de utilidad~(\ref{eq:utilidad}).
Se realizaron experimentos sistemáticos variando $\lambda \in \{0, 1, 5, 10\}$ y $\mu \in \{0{,}001, 0{,}01, 0{,}1\}$ en los entornos house y office con distintas densidades de actores.
Los resultados determinaron que los valores $\lambda = 5$ y $\mu = 0{,}01$ maximizan la puntuación compuesta en el conjunto de configuraciones evaluadas; estos valores se emplean de forma fija en todas las etapas de validación.
--- Completar con el análisis preliminar de calibración cuando esté disponible.%
]}

%%%%%%%%%%%%%%%%%%%%%%%%%%%%%%%%%%%%%%%%%%%%%%%%%%%%%%%%%%%%%%%%%%%%%%%%%%%%%%%%
